% =====================================================================
%  DANH MỤC THỰC NGHIỆM — trả lời câu hỏi "em đã làm những thí nghiệm gì?"
%  Biên dịch: xelatex danh_muc_thi_nghiem.tex
% =====================================================================
\documentclass[11pt,a4paper]{article}

\usepackage{fontspec}
\setmainfont{Times New Roman}
\usepackage{amsmath}
\usepackage{amssymb}
\usepackage{graphicx}
\usepackage{geometry}
\geometry{margin=2.2cm}
\usepackage{xcolor}
\usepackage{enumitem}
\setlist{nosep,leftmargin=1.3em}
\usepackage{titlesec}
\titleformat{\section}{\normalfont\large\bfseries\color{blue!55!black}}{\thesection.}{0.5em}{}
\titlespacing{\section}{0pt}{12pt}{5pt}
\usepackage{longtable}
\usepackage{array}

\newcommand{\nhan}[1]{\textbf{#1}}

\title{\textbf{DANH MỤC THỰC NGHIỆM}\\ \large
Dự án Contagion --- lan truyền prompt injection trong mạng LLM agent}
\author{Người báo cáo: \underline{\hspace{5cm}} \quad Ngày: \underline{\hspace{3.5cm}}}
\date{}

\begin{document}
\maketitle
\thispagestyle{empty}
\vspace{-0.8cm}

\begin{center}
\fbox{\parbox{0.94\textwidth}{
\textbf{Cách dùng.} Tài liệu này trả lời câu hỏi ``em đã làm những thí nghiệm gì''.
Mục~1 là bảng tóm tắt 10 nhóm thực nghiệm mang nội dung bài báo --- dùng để nói nhanh.
Mục~2 đi vào từng nhóm. Mục~3 là giai đoạn thăm dò (không dùng số trong bài) để giải
thích nếu thầy/cô hỏi vì sao có nhiều thư mục kết quả. Mục~4 là các phân tích chạy lại
từ dữ liệu đã có, \emph{không} tốn thêm chi phí gọi model.
}}
\end{center}

% =====================================================================
\section{Bảng tóm tắt 10 nhóm thực nghiệm}

\small
\begin{longtable}{|c|>{\raggedright\arraybackslash}p{3.5cm}|>{\raggedright\arraybackslash}p{4.3cm}|>{\raggedright\arraybackslash}p{6.3cm}|}
\hline
\textbf{\#} & \textbf{Nhóm thí nghiệm} & \textbf{Quy mô} & \textbf{Kết quả chính}\\
\hline
\endhead
1 & Đo lan truyền trên 5 model (ô chuẩn) & 5 model $\times$ (40 trial tự nhiên $+$ 30 trial/cạnh) & ASR: Llama $0{,}800$ · DeepSeek $0{,}475$ · qwen $0{,}300$ · Nova $0{,}075$ · Claude $0{,}000$\\
\hline
2 & Tính chuyển được (transportability) & 2 model, giao thức \texttt{fresh-artifact} & Llama lệch \textbf{$3{,}75\times$} ở cạnh giữa; DeepSeek khớp \textbf{$2{,}2\%$}\\
\hline
3 & Đối chứng chính sách artefact & 2 model $\times$ 2 chế độ & DeepSeek: $p$ từ $\mathbf{0{,}0030}$ (sai) $\to$ $\mathbf{0{,}630}$ (đúng) --- tự rút lại kết quả\\
\hline
4 & Hình thức thông điệp & 40 trial tự nhiên $+$ 30 phát lại/arm/cạnh & cùng mục tiêu, cùng cách chấm: $0{,}067$ vs $\mathbf{0{,}967}$ (\textbf{14 lần})\\
\hline
5 & Biến hình tấn công $\times$ defense & 3 model $\times$ 2 defense $\times$ 3 kiểu $\times$ 30 mẫu & redact chặn literal $1{,}00\to 0{,}00$ nhưng tách chuỗi qua mặt\\
\hline
6 & Topology & 3 topology $\times$ $n=7$ & ASR $0{,}450$ / $0{,}725$ / $0{,}800$; $R_0$ $0{,}821$ / $0{,}420$ / $0{,}831$\\
\hline
7 & Đường cong chiều sâu & 3 model $\times$ (60 trial $+$ 30/cạnh) & độ dốc $\rho$: $+1{,}00$ · $+0{,}77$ · $-0{,}07$\\
\hline
8 & Mô hình có vòng lặp & 2 model $\times$ 40 trial $\times$ 5 vòng & dự đoán $\rho$ đúng \emph{cả hai chiều}: bền (Llama) và yếu dần (qwen)\\
\hline
9 & Ảnh hưởng lên công việc hợp lệ (utility) & 2 model $\times$ 3 defense $\times$ 20 cặp & redact: ASR $0{,}450\to 0$, giữ được $1{,}000$; paraphrase: giữ được $0{,}400$ (tệ hơn không làm gì)\\
\hline
10 & Độ nhạy \& kiểm định giả định thống kê & 8 seed $\times$ 20 trial/cạnh; $+$ dữ liệu tổng hợp & $\varphi=0{,}58$ (i.i.d. đạt); $\chi^2=13{,}40$, $p=0{,}0014$ (phụ thuộc ngữ cảnh)\\
\hline
\end{longtable}

\normalsize
\textbf{Tổng khối lượng:} $\approx$ 9.070 lượt gọi model (cận dưới) · 2.342 trial tự nhiên ·
46 thư mục kết quả (25 có cấu trúc chuẩn) · $\approx$ 4,3 giờ máy ghi lại được.

% =====================================================================
\section{Chi tiết từng nhóm}

\subsection{Nhóm 1 --- Đo lan truyền trên 5 model (ô chuẩn của bài)}
\nhan{Mục đích:} trả lời ``model nào dễ bị lây, và xác suất sống sót qua mỗi lần chuyển
giao là bao nhiêu''.\\
\nhan{Thiết lập:} chuỗi 4 agent, đầu vào ở \texttt{agent\_0}; 40 lần chạy tự nhiên để đo
tỉ lệ thành công đầu-cuối; 30 lần cho \emph{mỗi} cạnh theo giao thức có kiểm soát để đo
xác suất sống sót có điều kiện; 3 ngữ cảnh công việc hợp lệ xoay vòng.\\
\nhan{Kết quả:} Llama 3.3 70B $0{,}800$; DeepSeek V3.2 $0{,}475$; qwen2.5:7b $0{,}300$;
Nova Pro $0{,}075$; Claude Sonnet 4.5 $0{,}000$. Xác suất sống sót chênh nhau
\textbf{6,5 lần} giữa các vai trò agent (worker $0{,}133$ so với reviewer $0{,}867$).\\
\nhan{Ý nghĩa:} ``dễ bị lây'' \emph{không} đi theo ``model mạnh hơn'' --- phải đo, không
suy từ bảng xếp hạng.

\subsection{Nhóm 2 --- Tính chuyển được: con số đo cách ly có dùng được không?}
\nhan{Mục đích:} đây là \textbf{câu hỏi trung tâm} của dự án. Các benchmark hiện nay đo
một agent rồi người ta nhân các con số đó lên để dự đoán cho pipeline --- phép nhân đó có
đúng không?\\
\nhan{Thiết lập:} chạy thêm giao thức \texttt{fresh-artifact} (mỗi trial rút một artefact
mới) và đo đồng thời xác suất sống sót \emph{trong chuỗi thật}, để so trực tiếp hai con
số.\\
\nhan{Kết quả:} DeepSeek: tích cách ly $0{,}409$ so với ASR $0{,}400$ --- khớp tới
\textbf{2,2\%}. Llama: cạnh giữa đo cách ly $0{,}233$ nhưng trong chuỗi thật $0{,}875$ ---
\textbf{thấp hơn $3{,}75$ lần}.\\
\nhan{Ý nghĩa:} giả định ngầm của mọi benchmark một-agent \emph{thất bại trên một số
model}, và mức thất bại phụ thuộc model.

\subsection{Nhóm 3 --- Đối chứng chính sách artefact (một kết quả của chính em đã bị rút lại)}
\nhan{Mục đích:} kiểm xem kết quả có phải do \emph{cách đo} mà ra không.\\
\nhan{Thiết lập:} chạy cùng một ô hai lần: một lần \emph{giữ cố định} artefact đã bị nhiễm
dùng lại cho mọi trial, một lần rút artefact mới cho từng trial.\\
\nhan{Kết quả:} với chế độ giữ cố định, DeepSeek cho $p=\mathbf{0{,}0030}$ --- nhìn như
bác bỏ giả thuyết Markov. Khi rút artefact mới: $p=\mathbf{0{,}630}$, \emph{không} bác bỏ
được gì. Ngược lại, bất thường của Llama \emph{đứng vững} qua cả hai chế độ.\\
\nhan{Ý nghĩa:} em giữ cả cái sai và cách sửa trong bài, vì với nghiên cứu đo lường thì
\textbf{cách đo là một phần của kết quả}.

\subsection{Nhóm 4 --- Hình thức thông điệp quyết định sự sống sót}
\nhan{Mục đích:} tìm \emph{cơ chế} giải thích vì sao phép đo cách ly sai.\\
\nhan{Thiết lập:} giữ nguyên mục tiêu, agent nhận, cách chấm và phân công; chỉ đổi cách
trình bày thông điệp: một câu khẳng định trần, so với nhúng trong một sản phẩm công việc.\\
\nhan{Kết quả:} $0{,}067$ so với $\mathbf{0{,}967}$ --- chênh \textbf{14 lần}. Tự kiểm tra
lại bằng cách phát lại nội dung bắt được: khớp \textbf{2/3} cạnh (báo cáo đúng 2/3, không
làm tròn thành ``khớp hoàn toàn'').\\
\nhan{Ý nghĩa:} artefact ``chuẩn hoá'' mà phép đo cách ly dùng \emph{không đại diện} cho
thông điệp thật.

\subsection{Nhóm 5 --- Biến hình tấn công $\times$ cơ chế phòng thủ}
\nhan{Mục đích:} đo xem một cơ chế phòng thủ có thật sự tác dụng, và nó bị lách thế nào.\\
\nhan{Thiết lập:} 3 model $\times$ 2 defense (không / redaction) $\times$ 3 kiểu trình bày
(thường, chèn khoảng trắng, tách chuỗi), 30 mẫu mỗi ô.\\
\nhan{Kết quả:} redaction chặn sạch dạng literal ($1{,}00 \to 0{,}00$) nhưng dạng tách
chuỗi lách qua được (tới $1{,}00$ trên DeepSeek). Trên Claude, khi \emph{không} phòng thủ
tỉ lệ đã chỉ $0{,}10$ --- nên ``redaction giảm về 0'' ở đó không chứng minh được gì.\\
\nhan{Ý nghĩa:} đây là \textbf{hiệu ứng sàn}: thiếu baseline ``không phòng thủ'' thì
không xác định được hiệu quả của defense.

\subsection{Nhóm 6 --- Topology: hình dạng mạng thay đổi cách rủi ro thể hiện}
\nhan{Thiết lập:} 3 topology cùng 7 agent (chuỗi, hình sao fan-in, hình cây), mỗi topology
40 trial tự nhiên $+$ 30 trial/cạnh.\\
\nhan{Kết quả:} ba topology có xác suất sống sót trung bình gần như nhau
($\bar s \approx 0{,}76$) nhưng: ASR $0{,}450$ / $0{,}725$ / $0{,}800$ và
$R_0$ $0{,}821$ / $0{,}420$ / $0{,}831$.\\
\nhan{Ý nghĩa:} \textbf{độ lan tới} và \textbf{khả năng tái lây} là hai rủi ro \emph{khác
nhau}; một topology có thể làm chúng tách rời. Hình sao tới đích nhiều hơn nhưng tái lây
ít hơn.

\subsection{Nhóm 7 --- Đường cong chiều sâu: sai số có tích luỹ không?}
\nhan{Thiết lập:} chạy \emph{một} chuỗi dài, mỗi nút là một mốc độ sâu; so tích các con số
đo cách ly với xác suất bị nhiễm thật ở độ sâu đó. Llama 15 agent, DeepSeek 8, qwen 7;
mỗi model 60 trial tự nhiên $+$ 30 trial/cạnh.\\
\nhan{Kết quả:} qwen sai số tăng \emph{đơn điệu} ($\rho=+1{,}00$, $+11{,}8$ điểm \%/hop);
Llama tăng ($+0{,}77$, $+6{,}8$); DeepSeek \emph{phẳng} ($-0{,}07$, $+0{,}1$). Bốn điểm
cuối của Llama bị \emph{loại} vì chuỗi chết hẳn, không gán sai số 100\%.\\
\nhan{Ý nghĩa:} \textbf{độ dốc của đường cong tự nó là phép chẩn đoán}: một chuỗi dài,
một lần chạy, biết ngay model đó có cộng dồn được các con số đo cách ly hay không.

\subsection{Nhóm 8 --- Mô hình có vòng lặp: tiêu chí ngưỡng và quy tắc thiết kế}
\nhan{Mục đích:} tiêu chí ``$R_0<1$ là an toàn'' vô hiệu trên mạng không có vòng (em
chứng minh bằng toán). Vậy khi mạng \emph{có} vòng thì sao?\\
\nhan{Thiết lập:} mô hình thật trong các framework: một manager giao việc cho 2 worker,
worker báo cáo ngược lại. Hai model (Llama, qwen), 40 trial $\times$ 5 vòng, cộng 30
trial/cạnh để đo các xác suất sống sót.\\
\nhan{Kết quả:} công thức $\rho=\sqrt{\sum_j a_jc_j}$ (kiểm bằng số tới
$4{,}4\times10^{-16}$ trên 320 cấu hình ngẫu nhiên) đưa ra \textbf{hai dự đoán ngược nhau}
và \emph{cả hai đều đúng}: Llama $\rho=1{,}342>1$, manager bị tái nhiễm $0{,}725$, tới
vòng 5 vẫn $0{,}700$ (bền); qwen $\rho=0{,}760<1$, manager từ $0{,}775$ giảm còn
$\mathbf{0{,}650}$ (yếu dần). Từ đó ra \textbf{quy tắc thiết kế}: vượt ngưỡng khi
$\sum_j a_jc_j>1$; đối xứng thì $w\,s^2>1$ --- ở $s=0{,}5$ cần 5 worker, $0{,}7$ cần 3,
$0{,}9$ cần 2.\\
\nhan{Giới hạn đã ghi trong bài:} mới 5 vòng --- đủ phân biệt ``yếu dần'' với ``bão hoà'',
\emph{chưa} đủ để khẳng định tắt hẳn.

\subsection{Nhóm 9 --- Ảnh hưởng lên công việc hợp lệ (utility)}
\nhan{Thiết lập:} mỗi cơ chế phòng thủ chạy 20 cặp: một lần sạch, một lần bị tấn công; đo
mạng còn hoàn thành đúng công việc hợp lệ được bao nhiêu phần.\\
\nhan{Kết quả:} trên DeepSeek, redaction chặn sạch tấn công ($0{,}450 \to 0$) mà
\emph{không} mất công việc hợp lệ (giữ được $1{,}000$); cơ chế viết lại ngữ nghĩa
(paraphrase) \emph{không giảm} lan truyền mà còn làm mất thêm ($0{,}400 < 0{,}600$ của
không làm gì).\\
\nhan{Ý nghĩa:} phải báo cáo \emph{cặp} (tỉ lệ tấn công, mức mất công việc hợp lệ); nếu
chỉ báo tỉ lệ tấn công thì paraphrase trông ``trung tính''.
\nhan{Hạn chế:} đại lượng này đang là \emph{proxy} (mức nhiễm bẩn sản phẩm), chưa phải độ
đúng thật.

\subsection{Nhóm 10 --- Độ nhạy và kiểm định giả định thống kê}
\nhan{Mục đích:} khoảng tin cậy dựa trên giả định các lần chạy độc lập --- giả định đó có
đúng với LLM không?\\
\nhan{Thiết lập:} 6 seed $\times$ 3 temperature, và một lần chạy riêng 8 seed ở một
temperature; cộng thêm kiểm định công cụ thống kê trên dữ liệu tổng hợp có đáp án biết
trước (60 lần lặp).\\
\nhan{Kết quả:} hệ số phân tán $\varphi=0{,}58$ trong cùng một temperature $\Rightarrow$
giả định độc lập \emph{được ủng hộ}. Nhưng kiểm định $\chi^2$ theo 3 ngữ cảnh công việc
\emph{bác bỏ} tính đồng nhất trên đúng cạnh yếu ($\chi^2=13{,}40$, $df=2$, $p=0{,}0014$;
$18/56$ so với $3/56$ so với $13/48$).\\
\nhan{Ý nghĩa:} xác suất sống sót \emph{phụ thuộc nội dung công việc hợp lệ} --- cùng cơ
chế mà nhóm 4 chỉ ra từ hướng khác. Các con số trong bài là trung bình trên ngữ cảnh.
Đây cũng là chỗ em báo cáo thẳng một kết quả \emph{bất lợi} cho giả định của chính mình.

% =====================================================================
\section{Giai đoạn thăm dò (không dùng số trong bài)}

Trước khi chốt thiết kế, em chạy 12 thư mục thí nghiệm nhỏ để trả lời các câu hỏi kỹ thuật.
Số liệu của chúng \textbf{không} xuất hiện trong bài báo --- mục đích chỉ là chốt cách đo:

\begin{sloppypar}
\begin{itemize}
  \item \textbf{Chấm điểm:} \texttt{asv\_dist\_probe}, \texttt{judge\_calib},
        \texttt{judge\_calib\_7b}, \texttt{judge\_calib\_smoke}, \texttt{judge\_calib\_smoke2}
        --- chọn ngưỡng và kiểm phân bố của công cụ chấm.
  \item \textbf{Chốt họ nhiệm vụ:} \texttt{mini\_study\_7b}, \texttt{mini\_study\_7b\_ctx},
        \texttt{mini\_study\_7b\_v2}, \texttt{task\_b\_pilot}, \texttt{semantic\_probe}
        --- thử nhiệm vụ ``đọc lại chuỗi'' (bão hoà ở $1{,}0$, không phân biệt được
        defense) rồi chuyển sang nhiệm vụ ``cạnh tranh ngữ nghĩa''.
  \item \textbf{Thăm dò defense:} \texttt{defense\_diag}, \texttt{defense\_reinj\_probe},
        \texttt{taskb\_obfuscation}, \texttt{task\_b\_redact\_nlarge},
        \texttt{rerun\_chain\_n40}, \texttt{rerun\_chain\_paraphrase\_n40}.
  \item \textbf{Kiểm kết nối:} \texttt{smoke\_bedrock} (2 trial, chỉ để chắc API chạy được).
\end{itemize}
\end{sloppypar}

\nhan{Vì sao phải nói phần này:} nếu chỉ đưa ra con số ``46 thư mục'', người nghe dễ hiểu
nhầm là 46 thí nghiệm đều dùng trong bài. Thực tế \textbf{25 thư mục} có cấu trúc chuẩn
được dùng, phần còn lại là chạy thử để chốt thiết kế. Một bài đo lường tử tế thì phải có
giai đoạn này, và nó cũng là thứ giúp phát hiện ra hai lỗi thiết kế lớn (xem nhóm 3).

% =====================================================================
\section{Các phân tích chạy lại từ dữ liệu đã có (0 lần gọi model)}

Đây là điểm em muốn thầy/cô lưu ý: \textbf{không phải việc gì cũng cần gọi model}. Một
phần đáng kể đóng góp của bài là phân tích lại dữ liệu đã đo, không tốn thêm chi phí:

\begin{itemize}
  \item \texttt{threshold\_analysis.py} --- tính bán kính phổ, kết luận về tiêu chí ngưỡng,
        và bài toán đặt phòng thủ ở đâu (Nhóm 6, §3.8 của bài).
  \item \texttt{isolation\_validity.py} --- đối chiếu xác suất sống sót đo cách ly với
        trong chuỗi (Nhóm 2).
  \item \texttt{markov\_formal\_all.py} --- kiểm định composition có khoảng tin cậy và
        độ lệch nhỏ nhất phát hiện được.
  \item \texttt{depth\_trend.py} --- thống kê hình dạng đường cong chiều sâu: tương quan
        hạng và độ dốc, có loại tường minh các điểm không xác định (Nhóm 7).
  \item \texttt{cyclic\_design\_rule.py} --- kiểm chứng công thức
        $\rho=\sqrt{\sum_j a_jc_j}$ và quy tắc thiết kế (Nhóm 8).
  \item \texttt{utility\_from\_outputs.py} --- chấm lại phần ảnh hưởng lên công việc hợp lệ
        từ đầu ra thô, \emph{không} cần gọi lại model (Nhóm 9).
  \item \texttt{audit\_numbers.py}, \texttt{scale\_report.py},
        \texttt{check\_results\_complete.py} --- tự kiểm: số trong bài có khớp file kết quả
        không, quy mô thực nghiệm, và có mất kết quả nào không.
\end{itemize}

\vspace{0.4cm}
\begin{center}
\fbox{\parbox{0.94\textwidth}{
\textbf{Nếu chỉ có 3 phút để nói về thực nghiệm:} em đã chạy
\textbf{10 nhóm thực nghiệm} trên \textbf{5 model / 5 nhà phát triển}, tổng
$\approx$ \textbf{9.070 lượt gọi model} và \textbf{2.342 trial tự nhiên}; trong đó có
\textbf{1 nhóm đối chứng} cho thấy một kết quả ban đầu của chính em là lỗi thiết kế đo và
đã bị rút lại. Ngoài ra có \textbf{7 phân tích} chạy lại từ dữ liệu đã có mà không tốn
thêm chi phí gọi model.
}}
\end{center}

% =====================================================================
\section{Lời nói gợi ý --- nói gì với giảng viên}

\subsection*{Nếu chỉ được nói 30 giây}
``Thưa thầy/cô, em làm thực nghiệm theo \emph{hai giao thức độc lập}: một giao thức đo
xác suất lây qua \emph{từng cạnh} một cách có kiểm soát, và một giao thức chạy cả mạng
\emph{tự nhiên} để đo tỉ lệ thành công đầu-cuối. Em chạy trên \textbf{5 model của 5 nhà
phát triển}, tổng khoảng \textbf{9 nghìn lượt gọi model}. Trong đó có một thí nghiệm đối
chứng cho thấy một kết quả ban đầu của chính em là do lỗi thiết kế đo, và em đã rút lại.''

\subsection*{Nếu được nói khoảng 3 phút --- nói theo 10 câu này, mỗi câu một thí nghiệm}

\begin{enumerate}
  \item \nhan{Đo lan truyền cơ bản (5 model):} mỗi model em chạy 40 lần tự nhiên và 30
        lần cho mỗi cạnh. Kết quả: Llama dễ lây nhất ($0{,}8$), rồi DeepSeek ($0{,}475$),
        qwen ($0{,}3$), Nova ($0{,}075$), Claude gần như không lây ($0$). Đáng chú ý là
        thứ tự này \emph{không} theo độ mạnh của model.
  \item \nhan{Thí nghiệm quan trọng nhất --- con số đo riêng lẻ có dùng được không:} em
        kiểm xem xác suất đo \emph{cách ly} từng cạnh có bằng xác suất \emph{trong chuỗi
        thật} không. Trên DeepSeek thì khớp (lệch $2{,}2\%$); trên Llama thì lệch
        \textbf{3,75 lần}. Nghĩa là cách làm quen thuộc ``đo một agent rồi nhân lên'' là
        \emph{không đáng tin}.
  \item \nhan{Thí nghiệm đối chứng về cách đo:} em chạy cùng một ô hai lần --- một lần
        dùng lại một mẫu cố định, một lần rút mẫu mới cho từng lần chạy. Lần đầu cho kết
        quả rất mạnh ($p = 0{,}003$), nhưng đó là \emph{lỗi do cách đo}; sửa lại thì
        $p = 0{,}63$, không kết luận được gì. Em rút lại kết quả đó và ghi cả cái sai lẫn
        cách sửa vào bài.
  \item \nhan{Tìm cơ chế --- vì sao sai:} giữ nguyên mục tiêu, agent nhận và cách chấm,
        chỉ đổi \emph{cách trình bày} thông điệp. Một câu khẳng định trần lây được
        $0{,}067$; nhúng trong một sản phẩm công việc lây được $0{,}967$ --- chênh
        \textbf{14 lần}.
  \item \nhan{Thử cơ chế phòng thủ:} redaction xóa chuỗi bí mật thì chặn được $100\%$
        dạng tấn công thường, nhưng nếu kẻ tấn công \emph{tách chuỗi} ra thì lách qua được.
        Và trên model vốn đã kháng tốt thì ``chặn được'' không chứng minh được gì.
  \item \nhan{So ba cách bố trí mạng cùng 7 agent:} chuỗi, hình sao, hình cây. Xác suất
        lây trung bình gần như nhau nhưng kết cục khác hẳn: hình sao tới đích nhiều hơn
        nhưng \emph{tái lây ít hơn}.
  \item \nhan{Chạy một chuỗi dài để xem sai số có tích luỹ:} Llama 15 agent, DeepSeek 8,
        qwen 7. Trên qwen sai số tăng đều từng bước, trên Llama cũng tăng, còn DeepSeek thì
        \emph{phẳng}. Nên độ dốc của đường cong tự nó là cách chẩn đoán.
  \item \nhan{Làm mô hình có vòng:} một manager giao việc cho hai worker rồi worker báo
        cáo ngược lại. Em tính ra công thức ngưỡng, công thức đó dự đoán \emph{hai chiều
        ngược nhau} trên hai model --- và em chạy cả hai: Llama thì lây dai dẳng, qwen thì
        yếu dần. \emph{Cả hai đều đúng như dự đoán.}
  \item \nhan{Đo ảnh hưởng lên công việc hợp lệ:} redaction chặn sạch tấn công mà không
        làm mất việc; còn cơ chế viết lại ngữ nghĩa không giảm được lây mà còn làm mất
        thêm --- \emph{tệ hơn cả không làm gì}.
  \item \nhan{Kiểm giả định thống kê:} các lần chạy có độc lập không. Trong cùng một
        temperature thì đạt, nhưng khi kiểm theo \emph{nội dung công việc} thì phát hiện
        xác suất lây phụ thuộc vào việc agent đang làm.
\end{enumerate}

\subsection*{Câu kết}
``Tổng lại: khoảng \textbf{9 nghìn lượt gọi model}, \textbf{5 model của 5 nhà phát triển},
\textbf{10 nhóm thực nghiệm}, và một kết quả của chính em đã bị \emph{rút lại} sau khi em
phát hiện lỗi trong thiết kế đo.''

\subsection*{Ba điều KHÔNG nên nói (dễ bị bắt lỗi)}
\begin{itemize}
  \item \textbf{Đừng nói ``em làm 46 thí nghiệm''.} Con số 46 là \emph{số thư mục kết quả};
        chỉ \textbf{25} thư mục có cấu trúc dùng cho bài. Nói: ``23 thí nghiệm chính, phần
        còn lại là chạy thử để chốt thiết kế''.
  \item \textbf{Đừng nói ``em bác bỏ được giả thuyết Markov''.} Kết quả đó đã bị rút lại
        ($p$ từ $0{,}003$ thành $0{,}63$). Nói đúng: ``em tưởng bác bỏ được, nhưng kiểm lại
        thì đó là lỗi thiết kế đo, nên em rút lại''.
  \item \textbf{Đừng nói vòng lặp ``tắt hẳn''.} Mới chạy 5 vòng --- chỉ đủ để nói
        ``yếu dần''. Nói ``yếu dần'' hoặc ``có dấu hiệu tắt dần''.
\end{itemize}

\subsection*{Nếu thầy/cô hỏi thêm}
\begin{itemize}
  \item \textbf{``Sao biết kết quả không phải do lỗi code?''} --- Em có hai giao thức độc
        lập cho hai con số phải khớp nhau theo một đẳng thức đã chứng minh; công cụ thống
        kê được kiểm trên dữ liệu tổng hợp có đáp án; và có một script đối chiếu từng con số
        trong bài với file kết quả gốc.
  \item \textbf{``Cỡ mẫu 40 có đủ không?''} --- Ở $n = 40$ chỉ phát hiện được lệch lớn hơn
        $0{,}26$. Em công bố con số này trong bài và \emph{không} coi kết quả ``không bác
        bỏ'' là bằng chứng.
  \item \textbf{``Tốn bao nhiêu tiền / thời gian?''} --- Phần chạy trên model cloud dùng
        kinh phí của bộ môn; phần qwen chạy tại máy nên miễn phí. Thời gian máy ghi lại
        được $\approx 4{,}3$ giờ, nhưng đây là cận dưới.
  \item \textbf{``Bước tiếp theo là gì?''} --- Mở rộng sang mạng lưới và mạng tranh luận
        (có đường song song --- nơi bài toán đặt phòng thủ thành bài toán tối ưu thật sự),
        và xây họ nhiệm vụ có đáp án kiểm chứng được để đo ảnh hưởng lên công việc chính xác
        hơn.
\end{itemize}

\end{document}
